% TOP_PROBLEM: {"id":30007726,"problem_number":"LOCAL-30007726","title":"Steering patients toward higher-value care","category_id":21,"category":{"id":21,"name":"economics","display_name":"Economics","description":"Open economic theory statements, published research agendas, and proposed scoped specifications; question provenance and historical priority are qualified per record.","slug":"economics","order_index":21,"created_at":"2026-10-08T00:00:00Z"},"status":"provenance_pending","external_url":null,"legacy_ids":[],"published":false,"statement":"Test whether information, scheduling help, and travel assistance improve net patient welfare in elective joint replacement, and evaluate how market saturation changes the result.","economics":{"original_id":215,"field":"Health economics","kind":"empirical","formalization_basis":"proposed_specification","source_status":"not_verified","priority_status":"not_established","priority_note":"No readable primary passage explicitly stating this precise open question was verified. A supporting paper, abstract, or topic citation is insufficient to establish origin or unresolved status.","earliest_verified_reference":null,"current_reference":{"authors":["Martin Gaynor","Kate Ho","Robert J. Town"],"title":"The Industrial Organization of Health Care Markets","year":2014,"url":"https://kateho.scholar.princeton.edu/sites/g/files/toruqf4136/files/kateho/files/healthcare_io_3_24_14.pdf","doi":"10.3386/w19800","locator":"§7 Models of Health Plan Choice, p. 40; §8, p. 43; §9 Future Directions, pp. 47–48","evidence_summary":"The authors explicitly discuss unanswered questions about consumer inertia, adverse selection, provider incentives, and consumer search frictions."},"formalization_note":"Proposed scoped research specification. The original catalogue prompt was a synthesis, and no canonical conjecture or verified original formulation is claimed. The supporting source, when retained, establishes context or existing results rather than this exact open question."},"statusQualification":"Provenance pending: an explicit published statement of this exact open question was not verified. This is a catalogue-authored candidate specification, not a certified open conjecture. Proposed scoped research specification. The original catalogue prompt was a synthesis, and no canonical conjecture or verified original formulation is claimed. The supporting source, when retained, establishes context or existing results rather than this exact open question.","statusReviewedAt":"2026-10-08"}
% FIRST_POSTED: 2026-10-08
% ENTRY_KIND: statement_only
% !TeX program = lualatex
\documentclass[11pt]{article}
\usepackage{fontspec}
\usepackage[a4paper,margin=25mm]{geometry}
\usepackage{amsmath,amssymb,mathtools}
\usepackage{xurl}
\usepackage[hidelinks]{hyperref}
\newcommand{\sref}[2]{\hyperlink{source:#1:#2}{[S#2]}}
\setlength{\emergencystretch}{3em}
\begin{document}
% problemId:30007726
\section*{Steering patients toward higher-value care}\label{30007726}
\subsection{Definitions and mathematical statement}
Consider privately insured adults referred for elective joint replacement in one regional health-care market. Let \(z\in\{0,1\}\) indicate an offer combining provider quality information, scheduling assistance, and a stated travel subsidy. Let \(J_i(z)\) be the chosen provider, \(H_i(z)\) one-year health improvement, \(C_i(z)\) real treatment cost, and \(T_i(z)\) patient time and travel costs. For monetary health valuation \(\lambda>0\), define
\[\tau_W=E[\lambda\{H_i(1)-H_i(0)\}-\{C_i(1)-C_i(0)\}-\{T_i(1)-T_i(0)\}]-k,\]
where \(k\) is per-eligible-person program resource cost. Provider quality is estimated independently of treatment assignment, with risk adjustment specified before the intervention. Identify \(\tau_W\) through randomized offers and report uptake and provider-choice effects without excluding patients who decline assistance. At larger market saturation, allow queues, prices, and provider capacity to respond; a small trial need not identify these responses. Determine which component of the offer matters only when factorial variation or credible mechanism restrictions permit. The health-market review raises consumer-friction research questions, but the precise joint steering intervention and welfare target were not explicitly verified there.

\par\textbf{Formulation and scope.} Proposed scoped research specification. The original catalogue prompt was a synthesis, and no canonical conjecture or verified original formulation is claimed. The supporting source, when retained, establishes context or existing results rather than this exact open question.

\par\textbf{Open-question provenance.} An explicit published open-question statement was not verified. Sources below are supporting literature and must not be treated as a first listing.

\par\textbf{Historical priority.} No readable primary passage explicitly stating this precise open question was verified. A supporting paper, abstract, or topic citation is insufficient to establish origin or unresolved status.

\subsection{Short English statement}
Test whether information, scheduling help, and travel assistance improve net patient welfare in elective joint replacement, and evaluate how market saturation changes the result.

\subsection{Sources}
\begin{itemize}\raggedright
\item[\textnormal{[S1]}]\hypertarget{source:30007726:1}{} Martin Gaynor, Kate Ho, Robert J. Town. 2014. The Industrial Organization of Health Care Markets. §7 Models of Health Plan Choice, p. 40; §8, p. 43; §9 Future Directions, pp. 47–48.
\url{https://kateho.scholar.princeton.edu/sites/g/files/toruqf4136/files/kateho/files/healthcare_io_3_24_14.pdf}.
DOI: 10.3386/w19800.
{\small\textit{Supporting or current literature; see evidence qualification. The authors explicitly discuss unanswered questions about consumer inertia, adverse selection, provider incentives, and consumer search frictions.}}
\end{itemize}
\end{document}
